% GENERATED by bin/curate/gen_tutorials_guide.py -- DO NOT EDIT BY HAND.
% Assembled from each case's controlDict description + header comment.

\clearpage
\tutentry{unsteady01\_startup\_transient}{tutorials/unsteady/unsteady01\_startup\_transient}
\begin{tutscheme}
\begin{tikzpicture}[>={Stealth[length=2mm]}, every node/.style={font=\scriptsize}]
  \node[draw, rounded corners=2pt, align=center, minimum width=1.9cm, minimum height=0.85cm, fill=black!4] (reactor) at (0.0,-0.0) {\textbf{reactor}\\[-1pt]{\tiny dynamicCSTR}};
\end{tikzpicture}
\end{tutscheme}
\tutwhat{Open-loop CSTR START-UP transient with OpenFOAM-style real-time instant dirs: charged cold (320 K), no control, watch T(t) and the holdup approach the natural steady state by scrubbing 0/ 50/ 100/ ... s}
\begin{tutnarrative}
Process: A continuous CSTR (first-order compA $\to$ compB, mild exotherm) is CHARGED cold (320 K, pure compA) and then left to run with NO controller. The jacket sits at a fixed 320 K. From the cold start the reactor warms a little from the reaction exotherm, the feed/ outlet convection competes, and (T, composition) relax to the natural open-loop steady state over \textasciitilde{} several residence times.

What this tutorial shows: The OpenFOAM-style transient solution directories. With `solutionControl \{ write true; \}` in controlDict, the case root fills with REAL-TIME directories `0/ 50/ 100/ ... 1000/`. Each `<t>/internalState` is the reactor HOLDUP at that instant; each `<t>/streamFaces` is the instantaneous outlet face. Scrub the times and you watch T(t) and the inventory march to steady state --- the same way you scrub a CFD transient in ParaView.

To compare CLOSED-LOOP control, see ctrl01\_cstr\_temp\_control (a PID drives T to a 350 K setpoint) and ctrl05\_disturbance\_transient (a PID rejecting a feed-temperature disturbance), both also writing time dirs.

\end{tutnarrative}
\begin{tutgolden}
\goldrow{kpi}{reactor}{T\_final}{320.759374011}
\goldrow{kpi}{reactor}{n\_compA\_final}{0.0117988900072}
\goldrow{kpi}{reactor}{n\_compB\_final}{0.000201109992805}
\goldrow{kpi}{reactor}{t\_end}{1000}
\goldrow{kpi}{balance}{energy\_balance\_available}{0}
\goldrow{kpi}{balance}{mass\_closure\_rel}{1.77265609598e-14}
\goldrow{kpi}{balance}{mass\_kg\_final}{0.36}
\goldrow{kpi}{balance}{mass\_kg\_in\_cum}{1.5}
\goldrow{kpi}{balance}{mass\_kg\_initial}{0.36}
\goldrow{kpi}{balance}{mass\_kg\_out\_cum}{1.5}
\goldrow{kpi}{balance}{mass\_residual\_kg}{-2.65898414398e-14}
\goldrow{kpi}{balance}{material\_available}{1}
\goldrow{kpi}{reactor}{F\_in\_final}{5e-05}
\goldrow{kpi}{reactor}{z\_in\_compA\_final}{1}
\goldmore{1}
\end{tutgolden}

\clearpage
\tutentry{unsteady02\_tanks\_in\_series}{tutorials/unsteady/unsteady02\_tanks\_in\_series}
\begin{tutscheme}
\begin{tikzpicture}[>={Stealth[length=2mm]}, every node/.style={font=\scriptsize}]
  \node[draw, rounded corners=2pt, align=center, minimum width=1.9cm, minimum height=0.85cm, fill=black!4] (tank1) at (0.0,-0.0) {\textbf{tank1}\\[-1pt]{\tiny dynamicCSTR}};
  \node[draw, rounded corners=2pt, align=center, minimum width=1.9cm, minimum height=0.85cm, fill=black!4] (tank2) at (2.9,-0.0) {\textbf{tank2}\\[-1pt]{\tiny dynamicCSTR}};
  \node[draw, rounded corners=2pt, align=center, minimum width=1.9cm, minimum height=0.85cm, fill=black!4] (tank3) at (5.8,-0.0) {\textbf{tank3}\\[-1pt]{\tiny dynamicCSTR}};
  \draw[->] ($(tank1.west)+(-1.05,0.00)$) -- ($(tank1.west)+(0,0.00)$) node[at start, above, font=\tiny, anchor=south west] {feed};
  \draw[->] (tank1) -- (tank2) node[midway, above, font=\tiny] {s12};
  \draw[->] (tank2) -- (tank3) node[midway, above, font=\tiny] {s23};
  \draw[->] ($(tank3.east)+(0,0.00)$) -- ($(tank3.east)+(1.05,0.00)$) node[at end, above, font=\tiny, anchor=south east] {product};
\end{tikzpicture}
\end{tutscheme}
\tutwhat{Tanks-in-series RTD: three equal mixing tanks chained by the ctrl router; impulse in tank 1 $\to$ Erlang-3 response at tank 3 (tbar = 3 tau, sigma2/tbar2 = 1/3), Levenspiel closed forms as anchors.}
\begin{tutnarrative}
the SECOND HALF of the RTD lesson, unlocked by the ctrl router (ruled 2026-08-24): three EQUAL mixing tanks in series, each tau = n/F = 0.004/5e-5 = 80 s, chained by the same `in`/`outputs` topology grammar every steady flowsheet already uses. The router applies each connection ONE-STEP-EXPLICITLY: at every accepted driver step the downstream tank integrates on the upstream outlet at the step's START -- a transport delay of one driver step (1 s here, against tau = 80 s), announced by the [routing] lines at startup.

THE EXPERIMENT. The tracer impulse is an INITIAL-CONDITION substitution in tank 1 only (compB 0.0002 of the 0.004 kmol holdup -- substitution, not addition, so every tank holds exactly tau = 80 s and the closed forms are exact). An initial holdup m in tank 1 is mathematically a perfect impulse into tank 1 at t = 0, so tank 3's outlet is the Erlang-3 density:

E(t) = t\textasciicircum{}2 e\textasciicircum{}(-t/tau) / (2 tau\textasciicircum{}3) tbar = 3 tau = 240 s sigma2 = 3 tau\textasciicircum{}2 = 19200 s\textasciicircum{}2 (sigma2/tbar2 = 1/3 exactly) peak at 2 tau = 160 s (visible in trajectory.csv)

Levenspiel, Chemical Reaction Engineering, ch. 14 (the tanks-in-series model): the variance ratio 1/N is THE diagnostic the lesson teaches -- one tank gives 1, three give 1/3, a PFR gives 0. ctrl18 is the N = 1 control; this case is N = 3 on identical tanks.

WHAT THE ANCHORS TOLERATE, stated: the engine measures moments by trapezoid on the accepted 1 s grid and truncates at 1600 s (= 20 tau of one tank); both effects are far inside the anchor bands, and the router adds a one-step transport delay per connection (2 s total on tbar's 240 s -- inside the 1 \% band, and the honest price of an explicit scheme, stated rather than hidden).

\end{tutnarrative}
\begin{tutgolden}
\goldrow{kpi}{balance}{energy\_balance\_available}{0}
\goldrow{kpi}{balance}{mass\_closure\_rel}{2.15876699233e-16}
\goldrow{kpi}{balance}{mass\_kg\_final}{0.36}
\goldrow{kpi}{balance}{mass\_kg\_in\_cum}{7.2}
\goldrow{kpi}{balance}{mass\_kg\_initial}{0.36}
\goldrow{kpi}{balance}{mass\_kg\_out\_cum}{7.2}
\goldrow{kpi}{balance}{mass\_residual\_kg}{-1.55431223448e-15}
\goldrow{kpi}{balance}{material\_available}{1}
\goldrow{kpi}{rtd}{area\_kmol}{0.000198752512114}
\goldrow{kpi}{rtd}{sigma2\_s2}{19198.7859557}
\goldrow{kpi}{rtd}{tbar\_s}{241.502454857}
\goldrow{kpi}{rtd}{truncation\_final\_over\_peak}{1.54894525802e-06}
\goldrow{kpi}{tank1}{F\_in\_final}{5e-05}
\goldrow{kpi}{tank1}{T\_final}{330}
\goldmore{23}
\end{tutgolden}

